\documentclass[a4paper,12pt]{article}
\usepackage[utf8]{inputenc}
\usepackage[T1]{fontenc}
\usepackage[italian]{babel}
\usepackage{geometry}
\usepackage{setspace}
\usepackage{amsmath}
\usepackage{listings}
\lstset{basicstyle=\small\ttfamily, breaklines=true, breakatwhitespace=false, columns=flexible}
\geometry{margin=2.5cm}
\onehalfspacing

\title{Verifica scritta di informatica -- Classe 4 AIQ -- Compito A\\
	Soluzioni (senza \texttt{EXISTS}, \texttt{ANY})}
\author{Prof. Fedeli Massimo}
\date{}

\begin{document}
	\maketitle
	
	\section*{Soluzioni -- Parte A (Query SQL)}
	
	\begin{enumerate}
		
		\item \textbf{Noleggi del cliente con id\_cliente = 7}
		\begin{lstlisting}
			SELECT n.id_noleggio, m.targa, n.data_inizio, n.data_fine,
			c.nome, c.cognome
			FROM NOLEGGI n
			JOIN MOTO m    ON m.id_moto    = n.id_moto
			JOIN CLIENTI c ON c.id_cliente = n.id_cliente
			WHERE n.id_cliente = 7;
		\end{lstlisting}
		
		\item \textbf{Noleggi iniziati a febbraio 2026}
		\begin{lstlisting}
			SELECT c.cognome, c.nome, m.targa, n.data_inizio, n.data_fine
			FROM NOLEGGI n
			JOIN CLIENTI c ON c.id_cliente = n.id_cliente
			JOIN MOTO m    ON m.id_moto    = n.id_moto
			WHERE n.data_inizio >= '2026-02-01'
			AND n.data_inizio <  '2026-03-01';
		\end{lstlisting}
		
		\item \textbf{Moto mai noleggiate}
		\begin{lstlisting}
			SELECT m.id_moto, m.targa, m.marca, m.modello
			FROM MOTO m
			LEFT JOIN NOLEGGI n ON n.id_moto = m.id_moto
			WHERE n.id_noleggio IS NULL;
		\end{lstlisting}
		
		\item \textbf{Clienti che hanno noleggiato almeno una Touring}
		\begin{lstlisting}
			SELECT DISTINCT c.id_cliente, c.cognome, c.nome
			FROM CLIENTI c
			JOIN NOLEGGI n ON n.id_cliente = c.id_cliente
			JOIN MOTO m    ON m.id_moto    = n.id_moto
			WHERE m.categoria = 'Touring';
		\end{lstlisting}
		
		\item \textbf{Noleggi senza alcun pagamento con esito positivo}
		\begin{lstlisting}
			SELECT n.id_noleggio, n.id_cliente, n.id_moto,
			n.data_inizio, n.data_fine
			FROM NOLEGGI n
			LEFT JOIN PAGAMENTI p
			ON p.id_noleggio = n.id_noleggio
			AND p.esito = 'positivo'
			WHERE p.id_pagamento IS NULL;
		\end{lstlisting}
		
		\item \textbf{Numero totale di noleggi per ogni moto}
		\begin{lstlisting}
			SELECT m.id_moto, m.targa,
			COUNT(n.id_noleggio) AS totale_noleggi
			FROM MOTO m
			LEFT JOIN NOLEGGI n ON n.id_moto = m.id_moto
			GROUP BY m.id_moto, m.targa;
		\end{lstlisting}
		
		\item \textbf{Spesa totale per cliente (solo pagamenti con esito positivo)}
		\begin{lstlisting}
			SELECT c.id_cliente, c.cognome, c.nome,
			COALESCE(SUM(p.importo), 0) AS spesa_totale
			FROM CLIENTI c
			LEFT JOIN NOLEGGI n    ON n.id_cliente  = c.id_cliente
			LEFT JOIN PAGAMENTI p  ON p.id_noleggio = n.id_noleggio
			AND p.esito = 'positivo'
			GROUP BY c.id_cliente, c.cognome, c.nome;
		\end{lstlisting}
		
		\item \textbf{Moto con costo totale manutenzione superiore al costo medio}
		
		\textit{Costo intervento:}
		$\text{ore\_manodopera} \times \text{costo\_manodopera\_oraria}
		+ \sum(\text{quantita} \times \text{costo\_unitario})$
		
		\begin{lstlisting}
			SELECT t.id_moto, t.targa, t.costo_manutenzione
			FROM (
			SELECT m.id_moto, m.targa,
			COALESCE(SUM(i.ore_manodopera * i.costo_manodopera_oraria), 0)
			+ COALESCE(SUM(ir.quantita * r.costo_unitario), 0)
			AS costo_manutenzione
			FROM MOTO m
			LEFT JOIN INTERVENTI i        ON i.id_moto       = m.id_moto
			LEFT JOIN INTERVENTI_RICAMBI ir ON ir.id_intervento = i.id_intervento
			LEFT JOIN RICAMBI r           ON r.id_ricambio   = ir.id_ricambio
			GROUP BY m.id_moto, m.targa
			) AS t
			WHERE t.costo_manutenzione > (
			SELECT AVG(t2.costo_manutenzione)
			FROM (
			SELECT m2.id_moto,
			COALESCE(SUM(i2.ore_manodopera * i2.costo_manodopera_oraria), 0)
			+ COALESCE(SUM(ir2.quantita * r2.costo_unitario), 0)
			AS costo_manutenzione
			FROM MOTO m2
			LEFT JOIN INTERVENTI i2         ON i2.id_moto        = m2.id_moto
			LEFT JOIN INTERVENTI_RICAMBI ir2 ON ir2.id_intervento = i2.id_intervento
			LEFT JOIN RICAMBI r2            ON r2.id_ricambio    = ir2.id_ricambio
			GROUP BY m2.id_moto
			) AS t2
			);
		\end{lstlisting}
		
		\item \textbf{Clienti che hanno noleggiato ogni categoria di moto attiva}
		
		\textit{Idea:} contare le categorie distinte noleggiate per cliente e confrontarle con il totale delle categorie attive.
		
		\begin{lstlisting}
			SELECT c.id_cliente, c.cognome, c.nome
			FROM CLIENTI c
			JOIN NOLEGGI n ON n.id_cliente = c.id_cliente
			JOIN MOTO m    ON m.id_moto    = n.id_moto
			WHERE m.attiva = 1
			GROUP BY c.id_cliente, c.cognome, c.nome
			HAVING COUNT(DISTINCT m.categoria) = (
			SELECT COUNT(DISTINCT m2.categoria)
			FROM MOTO m2
			WHERE m2.attiva = 1
			);
		\end{lstlisting}
		
		\item \textbf{Per ciascun cliente, il noleggio pi\`u recente}
		
		\textit{Criterio:} data\_inizio massima; a parit\`a di data, id\_noleggio massimo.
		
		\begin{lstlisting}
			SELECT c.id_cliente, c.cognome, c.nome,
			n.id_noleggio, n.id_moto,
			n.data_inizio, n.data_fine
			FROM CLIENTI c
			JOIN (
			SELECT n1.id_cliente,
			MAX(n1.data_inizio) AS max_data_inizio
			FROM NOLEGGI n1
			GROUP BY n1.id_cliente
			) x ON x.id_cliente = c.id_cliente
			JOIN NOLEGGI n
			ON n.id_cliente  = x.id_cliente
			AND n.data_inizio = x.max_data_inizio
			JOIN (
			SELECT n2.id_cliente, n2.data_inizio,
			MAX(n2.id_noleggio) AS max_id_noleggio
			FROM NOLEGGI n2
			GROUP BY n2.id_cliente, n2.data_inizio
			) y ON  y.id_cliente      = n.id_cliente
			AND y.data_inizio     = n.data_inizio
			AND y.max_id_noleggio = n.id_noleggio;
		\end{lstlisting}
		
	\end{enumerate}
	
	\section*{Soluzioni -- Parte B (Algebra Relazionale)}
	
	\textit{Nota:} si usa l'algebra relazionale estesa con l'operatore di raggruppamento $\gamma$.
	
	\begin{enumerate}
		
		\item \textbf{Targhe delle moto attive di categoria ``Touring''}
		\[
		\pi_{\text{targa}}\!\left(
		\sigma_{\text{attiva}=1\;\wedge\;\text{categoria}=\textquoteleft\text{Touring}\textquoteright}
		(\text{MOTO})
		\right)
		\]
		
		\item \textbf{Cognome e nome dei clienti registrati nel 2026}
		\[
		\pi_{\text{cognome},\,\text{nome}}\!\left(
		\sigma_{\textquoteleft\text{2026-01-01}\textquoteright
			\le\,\text{data\_reg}\,<\,
			\textquoteleft\text{2027-01-01}\textquoteright}
		(\text{CLIENTI})
		\right)
		\]
		
		\item \textbf{Targhe delle moto che compaiono in almeno un noleggio}
		\[
		\pi_{\text{targa}}\!\left(
		\text{MOTO}
		\bowtie_{\text{MOTO.id\_moto}=\text{NOLEGGI.id\_moto}}
		\text{NOLEGGI}
		\right)
		\]

		
		\item \textbf{Coppie (cliente, targa) per i noleggi di febbraio 2026}
		\[
		\pi_{\text{CLIENTI.id\_cliente},\,\text{MOTO.targa}}\!\left(
		\sigma_{\textquoteleft\text{2026-02-01}\textquoteright
			\le\,\text{data\_inizio}\,<\,
			\textquoteleft\text{2026-03-01}\textquoteright}
		(\text{NOLEGGI})
		\bowtie\,\text{CLIENTI}
		\bowtie\,\text{MOTO}
		\right)
		\]
		

		
	\end{enumerate}
	
\end{document}